\documentclass[a4paper]{article}

% ---------- Encoding, lingua, font ----------
\usepackage[utf8]{inputenc}
\usepackage[T1]{fontenc}
\usepackage{textcomp}
\usepackage[italian]{babel}
\renewcommand{\familydefault}{\sfdefault}
\emergencystretch=1.5em

% ---------- Colori (UNA sola volta, prima di tikz/tcolorbox) ----------
\usepackage[dvipsnames,table]{xcolor}

% ---------- Matematica ----------
\usepackage{amsmath, amssymb, amsfonts, amsthm}
\usepackage{mathtools}
\DeclarePairedDelimiter\ceil{\lceil}{\rceil}
\DeclarePairedDelimiter\floor{\lfloor}{\rfloor}
\usepackage{cancel}
\usepackage{gensymb}

% ---------- Tabelle, figure, utilità ----------
\usepackage{booktabs, xltabular, tabularx, makecell}
\usepackage{subcaption}
\usepackage{float}
\usepackage{graphicx}
\usepackage{import}
\usepackage{pdfpages}
\usepackage{xifthen}
\pdfminorversion=7
\pdfsuppresswarningpagegroup=1
\graphicspath{{./figures/}}

\renewcommand{\arraystretch}{1.5}
\newcolumntype{Y}{>{\raggedright\arraybackslash}X}

% ---------- TikZ / pgfplots ----------
\usepackage{tikz}
\usetikzlibrary{shapes, arrows, automata, petri, decorations.markings,
                decorations.pathreplacing, positioning, calc, quotes}
\usepackage{circuitikz}
\usepackage{forest}
\usepackage{pgfplots}
\pgfplotsset{compat=1.18, width=10cm}

\tikzset{
    place/.style={circle, thick, draw=black, minimum size=6mm},
    state/.style={circle, thick, draw=black, fill=white, minimum size=6mm},
}

% ---------- Palette (solo quella nuova) ----------
\definecolor{thmcol}{HTML}{2563EB}   % blu     - teoremi
\definecolor{lemcol}{HTML}{7C3AED}   % viola   - lemmi
\definecolor{defcol}{HTML}{059669}   % verde   - definizioni
\definecolor{excol}{HTML}{EA580C}    % arancio - esempi
\definecolor{remcol}{HTML}{6B7280}   % grigio  - osservazioni

% ---------- Listings (riusa la palette nuova) ----------
\usepackage{listings}
\lstdefinestyle{mystyle}{
    backgroundcolor=\color{remcol!8},
    commentstyle=\color{defcol},
    keywordstyle=\color{thmcol},
    numberstyle=\tiny\color{remcol},
    stringstyle=\color{lemcol},
    basicstyle=\ttfamily\footnotesize,
    breakatwhitespace=false, breaklines=true, captionpos=b,
    keepspaces=true, numbers=left, numbersep=5pt,
    showspaces=false, showstringspaces=false, showtabs=false, tabsize=2
}
\lstdefinestyle{terminal}{
    backgroundcolor=\color{black},
    basicstyle=\ttfamily\color{green},
    frame=single, breaklines=true, showstringspaces=false
}
\lstset{style=mystyle}

% ---------- Teoremi, definizioni, esempi (tcolorbox) ----------
\usepackage[most]{tcolorbox}
\tcbuselibrary{theorems}

\tcbset{
  thmstyle/.style={
    enhanced, breakable, sharp corners,
    boxrule=0pt, frame hidden,
    borderline west={3pt}{0pt}{#1},
    colback=#1!6,
    colbacktitle=#1!14, coltitle=#1!60!black,
    fonttitle=\bfseries,
    titlerule=0pt,
    separator sign none,
    description delimiters parenthesis,
    before skip=12pt, after skip=12pt,
    left=8pt, right=8pt, top=4pt, bottom=4pt,
    overlay unbroken and first={\tikzset{external/export=false}},
    overlay middle and last={\tikzset{external/export=false}},
  }
}

\newtcbtheorem[number within=section]{theorem}{Teorema}{thmstyle=thmcol}{thm}
\newtcbtheorem[use counter from=theorem]{lemma}{Lemma}{thmstyle=lemcol}{lem}
\newtcbtheorem[use counter from=theorem]{definition}{Definizione}{thmstyle=defcol}{def}
\newtcbtheorem[use counter from=theorem]{example}{Esempio}{thmstyle=excol}{ex}
\newtcolorbox{remark}[1][]{thmstyle=remcol, title={Osservazione}, #1}

% ---------- Esercizi (mdframed) ----------
\usepackage{mdframed}
\theoremstyle{definition}
\newmdtheoremenv[
    linecolor=gray, leftmargin=0, rightmargin=0,
    innertopmargin=8pt, innerbottommargin=8pt
]{exercise}{Esercizio}[section]

% ---------- Figure svg ----------
\newcommand{\incfig}[1]{%
    \def\svgwidth{\columnwidth}
    \import{./figures/}{#1.pdf_tex}
}

% ---------- Hyperref (sempre per ultimo, prima di externalize) ----------
\usepackage{hyperref}
\hypersetup{colorlinks=false}

% ---------- Externalize ----------
\usetikzlibrary{external}
\tikzexternalize[prefix=figures/tikz/, optimize=false]
\tikzexternaldisable


\begin{document}
\begin{titlepage}
	\begin{center}
		\vspace*{1cm}

		\Huge
		\textbf{Planning \& Reinforcement Learning}

		\vspace{0.5cm}
		\LARGE
		UniVR - Dipartimento di Informatica

		\vspace{1.5cm}

		\textbf{Tommaso Toffali}

		\vfill


		\vspace{0.8cm}


		1° Semestre 2026/2027
    \end{center}
\end{titlepage}

\tableofcontents
\pagebreak
\part{Classical Planning}
\section{Introduction to Planning}
\subsection{Historical Context and Motivation}
Planning has been a central topic in artificial intelligence since the inception
of the field at the \textbf{Dartmouth Conference in 1956}, where the foundings
fathers of AI first gathered to define the scope and ambitions of this new discipline.
Planning represent one of the foundamental capabilities required for intelligence behavior:
the ability to reason about sequences of actions that achieve desired goals.
\vspace{1em}
\noindent
\newline
The motivation for studying automated planning are manifold:
\begin{itemize}
    \item \textbf{Autonomous agents}: Enabling robots, software agents, and autonomouse system 
        to make decision and act independently in complex environemnts
    \item \textbf{Resource optimization}: Finding optimal or near-optimal stragies for Resource
        allocation, scheduling, and logistics
    \item \textbf{Scientific applications}: Modeling and solving problems in domain such as space 
        exploration, manufacturing, and healtcare
    \item \textbf{Theoretical foundations}: Understanding the computional complexity and formal properties
        of reasoning about action and change
\end{itemize}
\subsection{The planning Cycle}
The complete planning process involves several interconnected pahses that form a continuos cycle:
\begin{enumerate}
    \item \textbf{Plan Generation}: Given a description of the current state, available actions, and desired
            goals, synthesize a sequence of actions (a plan) that transforms the inital state into a goal state.
            This phase involves:
            \begin{itemize}
                \item Analyzing the initial state to understand what is currently true 
                \item Identifying the gap between the current state and the goal state
                \item Searching throungh the space of possible action sequences 
                \item Evaluating candidate plans for correctness and optimality 
                \item Selecting the best plan according to specified criteria (e.g., shortes length, minimum cost, fastest execution)
            \end{itemize}
    \item \textbf{Plan Deployment}: Execute the generated plan in the actual environment, monitoring the execution
            to detect deviations from expected behavior. This phase includes:
            \begin{itemize}
                \item Translating abstract plan actions into concrete executable commands 
                \item Monitoring sensors and feedback to verify that actions have the expected effects
                \item Detecting discrepancies between predicted and actual states
                \item Maintaining a record of executed actions for debugging and learning
            \end{itemize}
    \item \textbf{Replanning}: When execution monitoring detects that the current plan is no longer valid
            (due to unexpected events, actions failures, or environmental changes), generate a new plan
            or repair the existing one. Replanning strategies includes:
            \begin{itemize}
                \item \textbf{Plan repair}: Modify the existing plan minimally to accomodate the new situation 
                \item \textbf{Replan from scratch}: Generate an entirely new plan from the current state
                \item \textbf{Contingency planning}: Use pre-computed alternative plans for anticipated failures
            \end{itemize}
\end{enumerate}
This cycle reflects the reality that planning systems must operate in dynamic, partially unpredictable environments where
initial plans may need adaptation. The cyle continues iteratively until the goal is achived or deemed unreachable.
\subsubsection{Challenges in Real-World Planning}
Real-World deployment of planning systems faces several challengs:
\begin{itemize}
    \item \textbf{Excution uncertainty}: Actions may fail or have unexpected outcomes
    \item \textbf{Incomplete information}: The planner may not have complete knowledge of the environment
    \item \textbf{Dynamic environments}: The world may change while the plan is beign executed
    \item \textbf{Computional constraints}: Planning must often occur in real-time with limited computational resoruces 
    \item \textbf{Plan quality vs. planning time}: Trade-off between finding optimal plans and responding quickly
\end{itemize}
These challenge motivate the development of robust planning algorihms that cna handle unvertainty, adapt to changes,
and operate efficiently under resource constraints.
\section{Formalization of the Planning Problem}
\subsection{State Transition Systems}
The mathematical foundation of planning rests on the concept of state transition system, which provides a formal model of 
how actions transform states.

\begin{definition}{State Transition System}
A state transition system is a tuple \( \Sigma = (S,A,\gamma) \) where:
\begin{itemize}
    \item \( S \) is a finite or countably infinite set of \textbf{states}
    \item \( A \) is a finite or countably infinite set of \textbf{actions}
    \item \( \gamma : S \times A \rightarrow S \) is a \textbf{state transition function} that maps a state and an action to a resulting state
\end{itemize}
\end{definition}
The transition function \( \gamma(s,a) = s' \) specifies that executing action \( a \) in state \( s \) result in state
\( s' \). This function encodes the \textbf{dynamics} of the domain - how the world changes in response to actions.
\subsubsection{Properties of State Transition System}
State transition systems can be characterized by serveral important properties:
\begin{itemize}
    \item \textbf{Determinism}: In a deterministic system, \( \gamma \) is a function - each state-action pair leads to 
        exactly one successor state. In non-deterministic systems, multiple outcomes are possible.
    \item \textbf{Reachability}: A state \( s' \) is reachable from state \( s \) if there exists a sequence of actions 
        that transform \( s \) into \( s' \). The reachable state space from an initial state is often much smaller thant the totale
        state sapce.
    \item \textbf{Reversibility}: An action is reversible if there exists another action that undoes its effects. Many real-world
        actions are irreversible (e.g., breaking an object)
    \item \textbf{State space structure}: The connectivity and topology of the sate space significanyly affect planning complexity. Highly connected spaces 
        may have many solution paths, while sparse spaces may have few or no solutions.
\end{itemize}
\subsection{Classical Planning: Foundamental Assumptions}
Classical planning makes several simplyfying assumptions that restrict the class of the problems
considered but enable efficient algorithmic solutions.
\begin{definition}{Classical Planning Framework}
The classical planning framework relies on six key assumptions:
\begin{enumerate}
    \item Finitely many states
    \item Finitely many actions
    \item Deterministic transition function
    \item Full observability
    \item Instantaneous actions
    \item No exogenous events
\end{enumerate}
\end{definition}
\begin{enumerate}
    \item \textbf{Finitely many states}: The state space \( S \) is finite, allowing exhaustive search techniques.
        This assumption ensures that the planning is decidable and that search alogirhms will terminate
        \vspace{1em}
        \noindent
        \newline
        \textit{Jusitfication}: While real-world domains may have infinite state spaces (e.g. continuos 
        variables), finite approximations are often sufficient for pratical purposes.
    \item \textbf{Finitely many actions}: The action space \( A \) is finite, ensuring decidability. Each state
        has a finite branching factor in the state space graph.
        \vspace{1em}
        \noindent
        \newline
        \textit{Jusification}: Even in complex domains, the number of distinct action type is 
        typically manageable, though the number of ground actions (instantieted with specific objects)
        may be LARGE
    \item \textbf{Deterministic transition function}: For every state \( s \) and action \( a \), there 
        is exactlu one resulting state \( \gamma(s,a) \). Actions have predictable, certain outcomes.
        \vspace{1em}
        \noindent   
        \newline
        \textit{Justification}: Determinism simplified reasoning about action effects. Non-deterministic 
        planning requires more complex formalism (e.g., Markov Decision Processes)
    \item \textbf{Full observability}: The planner has complete knowledge of the current state at all times.
        There is no uncertainty about which state the system occupies.
        \vspace{1em}
        \noindent
        \newline
        \textit{Justification}: Full observability eliminates the need for belief state reasoning and sensing actions.
        Partial observability requires more sophisticated planning approaches.
    \item \textbf{Instantaneous actions}: Actions have no duration, they occur instantaneously. Temporal 
        reasoning is not required.
        \vspace{1em}
        \noindent
        \newline
        \textit{Justification}: Ignoring action durations simplifies the planning model. Temporal planning
        extends classical planning to handle durations and concurrent actions.
    \item \textbf{No exogenous events}: Only the planning agent can change the world through deliberate actions. 
        The environment remains static unless acted upon.
        \vspace{1em}
        \noindent
        \newline
        \textit{Justification}: Exogenous events (e.g., other agents, natural processes) introduce additional complexity.
        Classical planning assumes a static world between actions.
\end{enumerate}
These assumptions, while restrective, capture a significant and important class of planning problems and provide a foundation
for understanding more complex, realistic planning scenarios.
\subsubsection{Relaxing Classical Assumptions}
Modern planning research has devoloped extensions that relax these assumptions:
\begin{itemize}
    \item \textbf{Probabilistic planning}: Handles non-deterministic actions with probability 
        distributions over outcomes
    \item \textbf{Conformant planning}: Plans under partial observability without sensing 
    \item \textbf{Contingent planning}: Generates conditional plans that include sensing actions 
    \item \textbf{Temporal planning}: Reasons about actions durations and temporal constraints
    \item \textbf{Multi-agent planning}: Coordinates plans among multiple agents
\end{itemize}
\subsection{The Planning Problem}
Given a state transition system \( \Sigma = (S,A,\gamma) \) we can now formally define
the planning problem.
\begin{definition}{Planning Problem}
A planning problem is a tuple \( P = (\Sigma, s_0, G) \) where:
\begin{itemize}
    \item \( \Sigma = (S,A,\gamma) \) is a state transition system
    \item \( s_0 \in S \) in the \textbf{initial state}
    \item \( G \subseteq S \) is a set of \textbf{goal states}
\end{itemize}
\end{definition}
Alternatively, goals can be specified as a \textbf{goal formula} \( g \), a logical 
expression that is satisfied by exactly those states in \( G \). This allows more compact and expressive goal 
specifications.
\subsubsection{Goal specification}
Goals can be specified in several ways:
\begin{itemize}
    \item \textbf{Explicit goal states}: Enumerate the set \( G \) of accetable final states.
        This is impratical for large state spaces.
    \item \textbf{Goal formula}: A logical formula \( g \) such that 
        \( G = {s \in S \;|\; s \models g}\). For example, in propositional 
        logic : \( g = At(Robot,RoomB)\wedge Clean(RoomB) \)
    \item \textbf{Goal conditions}: A set of propositions that must be true in the goal state.
        This is the most common approach in classical planning.
    \item \textbf{Utility functions}: In optimization settings, goals may be specified as utility 
        functions to maximize or cost functions to minimize.
\end{itemize}
\begin{definition}{Solution Plan}
A \textbf{solution} to a planning problem \( P = (\Sigma,s_0,G) \) is a 
sequence of actions \( \pi = (a_1,a_2,\ldots,a_n) \) such that:
\[
  s_n = \gamma(\gamma(\ldots \gamma(\gamma(s_0,a_1),a_2)\ldots, a_{n-1}),a_n)\in G
\]
That is, executing the sequence of actions starting from the initial state \( s_0 \) results in a goal state.
\end{definition}
We can also write this more compactly using the notation \( \gamma^*(s,\pi) \) to denote
the state reached by executing plan \( \pi \) from state s:
\[
  \gamma^*(s_0,\pi) \in G
\]
\subsubsection{Plan Quality}
Not all solution plans are equally desirable. Common quality metrics include:
\begin{itemize}
    \item \textbf{Plan length}: Number of actions in the plan. Shorter plans are often preferred.
    \item \textbf{Plan cost}: Sum of actions costs. Each action \( a \) may have an associated cost \( c(a) \).
    \item \textbf{Makespan}: Total execution time (relevant in temporal planning).
    \item \textbf{Resource consumption}: Amount of resources used during execution.
\end{itemize}
An \textbf{optimal plan} minimizes the chosen quality metric among all solutions plans.
\section{Representation of Planning Problems}
\subsection{The challenge of explicit representation}
While the state transition system formalism is mathematically elegant, it faces a critical pratical 
challenge: \textbf{the curse of dimensionality}. Real-world planning domains can have exponentially large sate
spaces. For example:
\begin{itemize}
    \item A domain with \( n \) boolean variables has \( 2^n \) possible states 
    \item A domain with \( k \) objects and \( m \) binary relations has up to \( 2^{mk^2} \) states
    \item Enumerating all states and transitions explicitly is infeansible for even moderately-sized problems
\end{itemize}
\begin{example}{Blocks World Complexity}
Consider a Blocks World domain with \( n \) blocks. The number of possible 
configurations grows super-exponentially with \( n \). For \( n = 10 \) blocks,
there are more than \( 10^13 \) possible configurations. Explicitly representing the transition 
function would require storing information about trillions of state-action pairs.
\end{example}
This motivates the need for \textbf{compact, structured representation} that exploit regularities 
in the domain to represent large state spaces and transition functions implicitly.
\subsection{Action Schemas: Structured Representation}
The classical approach to compact representation uses \textbf{action schemas} (also called action 
template or operators). An action schema is a prametrized description of a family of related actions.
\begin{definition}{Action Schema}
An action schema consists of:
\begin{itemize}
    \item \textbf{Name and parameters}: A symbolic name and typed parameters, e.g., \( Move(x,y) \).
    \item \textbf{Preconditions}: A logical formula Pre specifying when the action can be executed
    \item \textbf{Effects}: A logical formula \( Eff \) specifying how the sate changes when the action in executed
\end{itemize} 
\end{definition}
\textbf{Actions} are \textbf{ground instances} of action schemas, obtained by binding the parameters to specific objects in the 
domain. For example, the schema \( Move(x,y) \) might generate ground actions \( Move(RoomA,RoomB) \), \( Move(RoomB,RoomC) \), etc.
\subsubsection{Advantages of Action Schemas}
Action schemas provide severeal benefits:
\begin{itemize}
    \item \textbf{Compactness}: A single schema can represent exponentially many ground actions 
    \item \textbf{Generality}: Schemas capture the general structure of actions independent of specific objects 
    \item \textbf{Scalability}: Adding new objects to the domain automatically generates new appliable actions 
    \item \textbf{Knowledge reuse}: Schemas learned in one domain can be transferred to similar domains.
\end{itemize}
\subsection{State Representation: Propositional logic}
In classical planning, states are typically represented using \textbf{propositional logic}:
\begin{itemize}
    \item A state is a set of \textbf{atoms} (propositional variables) that are true in that state
    \item Atoms not in the set are assumed false (\textbf{closed-world assumption})
    \item Action preconditions and effects are logical formula over these atoms 
\end{itemize}
The representation allows:
\begin{itemize}
    \item Compact encoding of structured states
    \item Efficient reasoning about action applicability and effects 
    \item Use of logical inference techniques for plan generation
\end{itemize}
\subsubsection{State Update Semantics}
\begin{example}{STRIPS Assumption}
When an action is executed, the state is updated according to the action's effects:
\begin{itemize}
    \item \textbf{Add list}: Atoms that become true after the action 
    \item \textbf{Delete list}: Atoms that become false after the action
    \item \textbf{Frame axioms}: Atoms not mentioned in the effects remain unchanged
\end{itemize}
The \textbf{STRIPS assumption} (named after the Stanford Research Institute Problem Solver) states
that only atoms explicitly mentioned in the effects changem all other atoms persist unchanged.
This simplifies reasoning about action effects.
\end{example}
\subsection{Example: Blocks World}
Consider the classic \textbf{Blocks World} domain, one of the most studied domain in planning research
\subsubsection{Domain Description}
The Blocks World consists of:
\begin{itemize}
    \item A set of blocks that can be stacked on top of each other 
    \item A table with unlimited space
    \item A robot arm that can pick up and put down blocks
\end{itemize}
\begin{example}{Blocks World State Representation}
\textbf{State predicates}:
\begin{itemize}
    \item \( \texttt{On}(x,y) \): Block \( x \) is directly on top of blocks \( y \).
    \item \( \texttt{OnTable}(x) \): Block \( x \) is on the table 
    \item \( \texttt{Clear}(x) \): Block \( x \) has no blocks on top of it 
    \item \( \texttt{Holding}(x) \): The robot arm is holding block \( x \).
    \item \( \texttt{ArmEmpty} \): The robot arm is not holding anything
\end{itemize}
\end{example}
\subsubsection{Action Schemas}
\textbf{PickUp}(x): Pick up block \( x \) from the table 
\begin{align*}
    &\text{Parameters: } x \text{(block)} \\
    &\text{Precondition: } \texttt{Clear}(x) \wedge \texttt{OnTable}(x) \wedge \texttt{ArmEmpty} \\
    &\text{Effect: } \texttt{Holding}(x) \wedge \lnot \texttt{OnTable}(x) \wedge \lnot \texttt{Clear}(x) \wedge \lnot \texttt{ArmEmpty}
\end{align*}
\textbf{PutDown}(x): Put down block \( x \) on the table
\begin{align*}
    &\text{Parameters: } x \text{(block)} \\
    &\text{Preconditions: } \texttt{Holding}(x) \\
    &\text{Effect: } \texttt{OnTable}(x) \wedge \texttt{Clear}(x) \wedge \texttt{ArmEmpty} \wedge \lnot \texttt{Holding}(x)
\end{align*}
\textbf{Stack}(x, y): Stack block \( x \) on top of block \( y \).
\begin{align*}
    &\text{Parameters: } x,y \text{(blocks)} \\
    &\text{Preconditions: } \texttt{Holding}(x) \wedge \texttt{Clear}(x) \\
    &\text{Effect: } \texttt{On}(x,y) \wedge \texttt{Clear}(x) \wedge \texttt{ArmEmpty} \wedge \lnot \texttt{Holding}(x) \wedge \lnot \texttt{Clear}(y)
\end{align*}
\textbf{Unstack}(x, y): Remove block \( x \) from on top of block \( y \).
\begin{align*}
    &\text{Parameters: } x,y \text{(block)} \\
    &\text{Precondition: } \texttt{Holding}(x) \wedge \texttt{Clear}(x) \wedge \texttt{ArmEmpty} \\
    &\text{Effect: } \texttt{Holding}(x) \wedge \texttt{Clear}(y) \wedge \lnot \texttt{On}(x,y) \wedge \lnot \texttt{Clear}(x) \wedge \lnot \texttt{ArmEmpty}
\end{align*}

\subsubsection{Example Problem Instance}
\textbf{Initial state}: Blocks A, B, C are all on the table
\[
  s_0 = \left\{ \texttt{OnTable}(A), \texttt{OnTable}(B), \texttt{OnTable}(C), \texttt{Clear}(A), \texttt{Clear}(B), \texttt{Clear}(C), \texttt{ArmEmpty}\right\}
\]
\textbf{Goal}: Stack the blocks in the order A on B on C 
\[
  G = \left\{ \texttt{On}(A,B), \texttt{On}(B,C), \texttt{OnTable}(C)\right\}
\]
\textbf{Solution plan}:
\begin{enumerate}
    \item \texttt{PickUp}(B)
    \item \texttt{Stack}(B,C)
    \item \texttt{PickUp}(A)
    \item \texttt{Stack}(A,B)
\end{enumerate}

\part{Advanced Planning Representation}
\section{Reduction to Propositional Satisfiability}
\subsection{Motivation}
We continue our discussion of the formalization and representation of planning problems.
A fundamental technique in automated planning is to \textbf{reduce the planning problem to the 
stisfiability problem of propositional logic}. This reduction allows us to leverage powerful SAT
solvers to find plans efficently.
\begin{definition}{Reduction to SAT}
The key idea is to encode:
\begin{itemize}
    \item The initial state as a propositional formula
    \item The goal conditions as a propositional formula 
    \item The action effects and preconditions as propositional formulas 
    \item The constraint that actions must form a valid sequence
\end{itemize}
If the resulting formula is satisfiable, the satisfying assignment corresponds to a valid plan.
\end{definition}
\subsection{Example Domain: Dock Worker Robots (DWR)}
To illustrate these concepts, we introduce the \textbf{Dock Worker Robots (DWR)} domain, a classic 
planning benchmark that models container logistics at a port.
\subsubsection{Domain description}
The DWR world consists of:
\begin{itemize}
    \item \textbf{Docks}: Three docks \( d_1, d_2, d_3 \) arranged in a triangular configuration 
    \item \textbf{Connectivity}: Docks are connected as follows: \( d_1 \leftrightarrow d_3, d_1 \leftrightarrow d_2, d_2 \leftrightarrow d_3 \).
    \item \textbf{Piles}: Storage locations at docks where containers can be stacked 
    \item \textbf{Containers}: Object that need to be moved between locations 
    \item \textbf{Robots}: Mobile agents that can move containers between piles 
\end{itemize}
\begin{center}
    \begin{tikzpicture}[
        node distance=1.5cm,
        every node/.style={circle, draw, minimum size=1.2cm}
    ]
    \node (d1) {$d_1$};
    \node (d2) [right=of d1] {$d_2$};
    \node (d3) [below=of d1] {$d_3$};

    \draw[<->] (d1) -- (d2);
    \draw[<->] (d1) -- (d3);
    \draw[<->] (d2) -- (d3);
    \end{tikzpicture}
\end{center}

\subsection{Initial State configuration}
The initial state of our example problem is:
\begin{itemize}
    \item \textbf{At dock} \( d_1 \): Pile \( p_1 \) contains two containers: \( c_1 \) (on top) and \( c_2 \) (below)
    \item \textbf{At dock} \( d_2 \):
    \begin{itemize}
        \item Pile \( p_2 \) contains one container: \( c_3 \).
        \item Pile \( p_3 \) is empty 
    \end{itemize}
    \item \textbf{At dock} \( d_3 \): Pile \( p_4 \) is empty
\end{itemize}

RIPRENDERE APPUNTI FINO A SLIDE 46
\section{Conflict-Driven Clause Learning (CDCL)}
Modern SAT solvers use \textbf{Conflict -Driven Clause Learning (CDCL)}, an advanced 
technique that combines depth-first search with intelligence learning from conflicts 
to efficently solve large satisfiability problems.
\subsection{CDCL Algorithm Overview}
CDCL implements \textbf{depth-first search strategy} with intelligent backjumping (undo one or  more decision) 
and clause learning:
\begin{definition}{CDCL Components}
\centering \textbf{CDCL} combines several key techniques:
\begin{itemize}
    \item \textbf{Unit propagation}: Efficient inference of forced assignments
    \item \textbf{Conflict analysis}: Learning from failures to guide future search 
    \item \textbf{Backjumping}: Intelligent backtracking to erliaer decision levels
    \item \textbf{Clause learning}: Adding learned constraints to prevent repeated conflicts
\end{itemize}
\end{definition}
\subsection{Backjumping: Generalized Backtracking}
Backjumping is a generalization of the basic backtracking (undo one decison) mechanism that allows 
the solver to jump back to erlier decision levels, not just the immediately preceding one.
\begin{definition}{Backjumping vs Backtracking}
\centering In \textbf{backjumping}, when a conflict occurs, the solver 
analyzes the conflict cause to determine the earliest decisione level that contributed to conflict.
It then jumps back to that level, rather than just the immediately preceding decision.
\vspace{1em}
\noindent
\newline
\textbf{Advantage}: Backjumping can skip over irrelevant decision levels, significantly reducing 
the search space compared to basic backtracking.
\end{definition}
\subsubsection{Conflict analysis example}
COnsider a conflict clause \( C = L_1 \vee \ldots \vee L_k \) where all literals are currently false. The analysis 
determines:
\begin{itemize}
    \item Why each literal \( L_i \) is false in the current trail \( \varGamma \).
    \item The decision level at which each literal became false
    \item The erliest decision level that contributed to the conflict
\end{itemize}
The solver then backjumps to that earliest relevant decision level, effectively skipping decision 
levels that did not contribute to the conflict.
\section{Binary Resolution}
The theoretical foundation of conflict analysis relies on \textbf{binary resolution}, a foundamental inference rule in propositional logic
\subsection{Resolution Rule}
\begin{definition}{Binary Resolution}
\centering Given two clauses containing complementary literals:
\[
  \frac{C \vee L \quad D \vee \lnot L}{(C \vee D)}
\]
\end{definition}
The \textbf{resolvent} \( C \vee D \) is logically entailed by the two premises. This rule allows deriving new 
clauses that preserve satisfiability.
\subsection{Soundness of Resolution}
Binary resolution is \textbf{sound}, the resolvent is logically entailed by the permise.
\begin{theorem}{Resolution Soundness}
\centering \( If {C \vee L, D \vee \lnot L} \models R where R is any formula, then {C \vee
L, D \vee \lnot L} \models (C \vee D) \).
\vspace{1em}
\noindent
\newline
\textit{Proof}. Consider any interpretation \( v \):
\begin{itemize}
    \item \( If v \models L \), then since \( v \models D \vee \lnot L \), we must 
        have \( v \models D \). Since \( v \models C \vee L \), we have \( v \models C \vee D \).
    \item \( If v \models \lnot L \) then since \( v \models C \vee L \), we must have \( v \models C \). Since \( v \models D \vee \lnot L \),
        we have \( v \models C \vee D \).
\end{itemize}
In both cases, \( v \models C \vee D \), so the resolvent is logically entailed by the premises.
\end{theorem}
\subsection{Application in Conflict Analysis}
In conflict analysis, resolution is used to explain why conflicts occur and to determine which
decision levels are responsible:
\begin{itemize}
    \item When a literal \( L_k \) in the conflict clause is false due to propagation (not a decision)
    \item The propagation was justified by a clause \( D \vee \lnot L_k \).
    \item By resolving the conflict clause with the justification clause, we can trace back the decision that ultimately caused the conflict.
\end{itemize}
\section{Explanation Mechanism}
The explanation mechenism in CDCL traces the implications that led to a conflict,
allowing the solver to learn more general constraints.
\begin{theorem}{Explanation rule}
\centering \[ \langle S \Vert \varGamma \Vert C \rangle \rightarrow \langle S \Vert \varGamma \Vert C \vee D \rangle \]
if \( \lnot L \in \varGamma \) where \( \lnot L \) was propagated using the clause \( D \vee \lnot L \).
\end{theorem}
\subsection{Explanation Process}
The explanation mechanism works as follows:
\begin{enumerate}
    \item Start with a conflict clause \( C = L_1 \vee \ldots \vee L_k \).
    \item For each literal \( L_i \) that is false due to propagation (not a decision):
    \begin{itemize}
        \item Find the justification clause that forced \( \lnot L_i \) into the trail 
        \item Resolve the current conflict clause with the justification clause 
        \item Continue untile all literals in the conflict clause are explained by decision 
    \end{itemize}
    \item The final resolved clause contains only literals that were decided (not propagated)
    \item This clause can be learned and added to the formula to prevent similar conflicts
\end{enumerate}
\part{SAT-Based Techniques}
\section{Modern SAT Solving}
\subsection{CDCL: An evolution of DPLL}
CDCL (\textbf{Conflict-Driven Clause Learning}) represent an evolution of the DPLL decision procedure 
for SAT. While CDCL maintains the same basic decision rule as DPLL, it introduces sophisticated conflict anaysis and learning mechanisms that 
dramatically improve solver efficiency on real-world instances.
\subsubsection{CDCL Transition Rules}
The CDCL procedure employs the following transition rules in search mode:
\begin{definition}{Decision rule}
\[
  \langle S \Vert \varGamma \rangle \rightarrow \langle S \Vert \varGamma, L^d \rangle
\]
if \( L \notin \varGamma \) and \( \lnot L \notint \varGamma \), where:
\begin{itemize}
    \item \( S \) is the set of clauses 
    \item \( \vaGamma \) is the trail (sequence of assigned literals)
    \item \( L^d \) is a decision literal (chosen by heuristic)
\end{itemize}
\end{definition}
\begin{definition}{Propagation rule}
\[
  \langle S \Vert \varGamma \rangle \rightarrow \langle S \Vert \varGamma, L_{C \vee L} \rangle 
\]
if \( C \vee L \in S, L \notin \varGamma, \lnot L \notint \varGamma \), and \( \varGamma \models \lnot C \), where:
\begin{itemize}
    \item \( C \) is a disjunction of literals (all false under the current trail)
    \item \( L \) is the unit literal that must be propagated 
    \item The subscript \( C \vee L \) indicates the justification clause for the propagation
\end{itemize}
\end{definition}
\begin{definition}{Success rule}
\[
  \langle S \Vert \varGamma \rangle \rightarrow "sat"
\]
if \( \varGamma \models S \) (i.e., the trail satisfies all clauses in the set)
\end{definition}
\begin{definition}{Conflict rule}
\[
  \langle S \Vert \varGamma \rangle \rightarrow \langle S \Vert \varGamma \Vert C \rangle 
\]  
if \( C \in S \) and \( \varGamma \models \lnot C \) (i.e., the trail conflicts with clause \( C \)), where
\( C \) is the conflict clause that triggers transition to conflict analysis mode.
\end{definition}
This four rules define the behavior of CDCL in \textbf{search mode}, which is identical to the DPLL
procedure. The key difference lies in how CDCL handles conflicts through its conflict analysis 
and learning machanisms.
\subsection{Decision Levels in CDCL}
The trail \( \varGamma \) in CDCL is organized into \textbf{decision levels} that provides
structure ti the search process:
\begin{itemize}
    \item The trail begins at \textbf{decision level 0} (empty trail)
    \item Each decision literal opens a new decision level 
    \item If the maximum decision level in \( \varGamma \) is \( n \), the next decision opens level \( n + 1 \).
    \item All propagations following a decision remain at the same level untile the next decision
\end{itemize}
\subsubsection{Trail as a Stack Structure}
The trail can be conceptualized as a \textbf{stack structure}:
\begin{enumerate}
    \item Initially, the stack is empty (decision level 0)
    \item A decision literal is puhsed onto the stack, opening a new level
    \item Unit propagations following the decision are pushed onto the stack at the same level 
    \item Another decision is maade, opening the next level 
    \item The process continues with alternating decision and propagations
    \item During backtracking/backjumping, literals are popped from the stack until reaching the target decision level
\end{enumerate}
This organization enables efficient backjumping by allowing the solver to quickly identify and 
return to relevant decision levels during conflict analysis.

\subsection{Illustrative example}
\begin{example}{Instance}
\centering \( S = \left\{ P, \lnot P \vee Q \vee R, T \vee \lnot Q, \lnot R \vee Q \right\} \)
\end{example}
The inital propagation and decision sequence proceeds as follows·
\[
  \langle S \Vert \epsilon \rangle \xrightarrow{prop} \langle S \Vert P_P \rangle
\]
\end{document}